\documentclass[11pt]{article}
\usepackage{geometry}
\geometry{
	a4paper,
	margin=2cm,
	left=3cm}
\usepackage[utf8]{inputenc}
%\usepackage[german]{babel}
\usepackage[T1]{fontenc}
\usepackage{graphicx}

\setlength{\parindent}{0em} %setzt Paragraphen-Einzug auf 0

\title{Firmenporträt Kinderarztpraxis}

\begin{document}
\includegraphics[scale=0.9]{banner.pdf}

\section*{Kinderarztpraxis Dr. med. Lasse Niessen}
Spielerstr. 56\\
01069 Dresden\\\\

Im Jahr 1974 wurde die Kinderarztpraxis von Dr. med. Carola Garten gegründet. Während ihrer langjährigen Dienstzeit baute sie einen großen Kundenstamm auf, der bei der Praxisübernahme von Dr. med. Lasse Niessen 2005 übernommen wurde. Neben den bereits existierenden Kunden übernahm Dr. Niessen auch das Praxisteam, welches aus vier Arzthelfern/innen besteht.\\

Unsere Praxis befindet sich im Zentrum Dresdens und besitzt eine gute Anbindung an die öffentlichen Verkehrsmittel. Der Wartebereich ist für ca. 30 Personen ausgelegt und wird vor allem durch unseren Spielbereich für die wartenden Kinder geprägt. Neben dem Spiel- und Malbereich gibt es auch eine reichliche Auswahl an Fachzeitschriften zu aktuellen Themen.\\

2011 wurde die Praxis erweitert, sodass nun auch Röntgenaufnahmen sowie eine Untersuchung mittels Ultraschall möglich ist. Somit ist eine umfangreichere Diagnose unserer Patienten möglich.\\

Seit 2012 kooperieren wir mit dem Uniklinikum Dresden um einen aktuellen Stand bzw. genügend Parallelität zur modernen Medizin zu wahren. Damit können wir unseren Patienten einen guten Zugang zu weiteren Behandlungsmöglichkeiten gewähren.\\

%todo
Durch unsere hell und freundlich gestalteten Warte- und Praxisräume möchten wir eine angenehme und stressarme Atmosphäre schaffen.
Wir schaffen ein familienfreundliches Umfeld, indem wir einfühlsame Mitarbeiter/innen beschäftigen und die Wartezeit für unsere jungen Patienten und deren Eltern besonders angenehm gestalten wollen. Wir bieten darüber hinaus alternative Heilmethoden an.
Durch unseren großen Spielbereich sollten die Kinder ausreichend Beschäftigungsmöglichkeiten haben.

\vspace*{\fill}
\includegraphics[scale=0.9]{footer.pdf}
\thispagestyle{empty}
\end{document}